\chapter*{CONCLUSION GÉNÉRALE}
\addcontentsline{toc}{chapter}{CONCLUSION GÉNÉRALE}
La reconstruction d'images en tomodensitométrie constitue un enjeu central de l'imagerie médicale moderne, où la nécessité de concilier qualité diagnostique, rapidité d'acquisition et réduction de dose impose de repenser les approches classiques. Les méthodes analytiques, bien qu'efficaces dans des contextes idéaux, montrent leurs limites face aux contraintes réelles de bruit et de sous-échantillonnage. Dans ce contexte, le cadre du \textbf{Compressed Sensing} s'impose comme une alternative fondamentale en exploitant la parcimonie des images pour permettre une reconstruction fiable à partir de données incomplètes, transformant ainsi un problème inverse mal posé en une formulation mathématiquement maîtrisable.\\

L'évolution vers des méthodes itératives régularisées marque une transition décisive dans la manière d'aborder la reconstruction tomographique. En formulant le problème sous la forme d'un modèle inverse intégrant à la fois la fidélité aux données et des contraintes de régularité, ces approches permettent de mieux préserver les structures essentielles de l'image. Cette dynamique met en évidence l'importance des outils mathématiques avancés, notamment les représentations parcimonieuses et les techniques d'optimisation convexe, qui offrent un compromis efficace entre robustesse face au bruit et précision des détails reconstruits.\\

Dans ce travail, on a proposé une approche reposant sur la combinaison de l'algorithme ADMM et des ondelettes biorthogonales Bior4.4. Ce choix s'inscrit dans une volonté d'exploiter simultanément la capacité de parcimonie offerte par cette base d'ondelettes, particulièrement adaptée à la préservation des contours et des structures géométriques, et l'efficacité de l'ADMM dans la résolution de problèmes d'optimisation décomposés. Le schéma proposé permet ainsi de dissocier les contraintes de fidélité et de régularisation, tout en assurant une convergence stable et rapide grâce à des étapes de mise à jour simples et bien conditionnées.\\

Les résultats obtenus confirment la pertinence et l'efficacité de cette stratégie, avec des performances quantitatives atteignant un PSNR moyen de $26.20dB$ et un indice SSIM de $0.866$, traduisant une excellente fidélité structurelle des images reconstruites. Comparée aux méthodes de référence, notre approche permet une amélioration notable de la qualité visuelle tout en conservant les détails anatomiques essentiels. Elle démontre ainsi sa capacité à produire des images fiables à partir de données fortement sous-échantillonnées, constituant un levier important pour la réduction de dose en pratique clinique. Ces résultats ouvrent des perspectives prometteuses, notamment en termes d'optimisation computationnelle, d'extension à la reconstruction tridimensionnelle et d'intégration avec des approches hybrides combinant méthodes variationnelles et apprentissage profond.

